\chapter{\label{chap:systemArchitecture}System Architecture}

Twin2MultiCloud is implemented as a research proof of concept that turns a
typed Digital Twin workload into cost and deployment evidence for Amazon Web
Services (AWS), Microsoft Azure, and Google Cloud Platform (GCP). Its purpose
is not to provide a general cloud-management platform. Instead, the system
operationalizes one fixed architecture, compares only functionally admissible
provider assignments, and preserves the trace from input through destruction.

\section{Overview and Design Principles}
\label{sec:arch-overview}

The deployed architecture is identified as
\texttt{six-layer-eventing@1}. It contains the five responsibilities inherited
from the theoretical Digital Twin model---ingestion, processing, tiered
storage, Twin state, and visualization---and makes Eventing an independent
mandatory responsibility. Storage is represented by separate hot, cool, and
archive components. Consequently, the optimizer assigns eight logical
components while the architecture remains a Six-layer model.

Five principles constrain the implementation:

\begin{enumerate}
    \item \textbf{Closed-world execution.} One versioned profile, component
    catalog, workload contract, and deployment contract define everything that
    can be executed. Runtime profile registration and arbitrary topology
    synthesis are outside the proof of concept.
    \item \textbf{Functional admissibility before price.} A provider
    assignment is costed only after every required capability, component,
    edge, runtime artifact, region binding, and cleanup contract can be
    resolved.
    \item \textbf{Immutable evidence.} The selected calculation, architecture,
    deployment specification, manifest, and provider-account bindings are
    digest-linked. A deployed Twin is not silently recalculated or updated in
    place.
    \item \textbf{Explicit cost-incurring actions.} Provider preparation,
    Apply, and Destroy require distinct review and confirmation boundaries.
    Reconnecting a user interface must not repeat any of these actions.
    \item \textbf{PoC-sized ownership.} One local owner profile, typed Twin
    interchange, bounded user functions, and durable operations are retained.
    Multi-tenancy, external application login, dashboard administration, and
    arbitrary deployment projects are not implemented.
\end{enumerate}

The system comprises a Flutter client and three Python services. The Flutter
client communicates only with the Management API. The Management API owns
users, Twins, CloudConnections, lifecycle state, confirmations, and evidence.
The Cost Optimizer owns capability resolution and cost calculation. The Cloud
Deployer owns deployment-package validation, Terraform execution, runtime
verification, and cleanup.

\section{Service and Trust Boundaries}
\label{sec:microservices}

\subsection{Service Communication}
\label{sec:comm-patterns}

All user-facing requests enter through the Management API. Calls to the
Optimizer and Deployer are server-to-server operations with explicit timeout
and error translation. Cloud credentials therefore never pass through the
Flutter client after storage, and the client cannot bypass Management-owned
immutability or confirmation checks.

The main cross-service evidence chain is:

\begin{center}
\texttt{Twin intent} $\rightarrow$ \texttt{CalculationRun}
$\rightarrow$ \texttt{RTA v2} $\rightarrow$ \texttt{RDS v2}
$\rightarrow$ \texttt{Manifest v4} $\rightarrow$
\texttt{DeploymentOperation}.
\end{center}

The Resolved Twin Architecture (RTA) records the selected logical component
assignments, resolved edges, workload reference, pricing evidence, and cost
totals. The Resolved Deployment Specification (RDS) expands these logical
decisions into exact implementation components, dimensions, identity
requirements, package artifacts, Terraform bindings, verification steps, and
cleanup requirements. Deployment Manifest v4 binds both documents, their
digests, the operation package, and the selected CloudConnections. Any digest
or version mismatch fails before Terraform execution.

\subsection{Credential Security and User Boundary}
\label{sec:credential-security}

The PoC retains one local owner profile because Twins and CloudConnections
need an explicit ownership boundary. Startup binds a configured static local
bearer to that profile. Google OAuth, Microsoft application login, university
SAML, JWT issuance, roles, and multi-tenant session management are not on the
runtime path. The existing login screen remains compiled as a dormant adapter
boundary but has no route.

Each provider credential is stored as an encrypted, user-owned
\emph{CloudConnection}. A user may retain several connections per provider and
select exactly one for every provider required by a resolved graph. Responses
contain metadata and validation state, never credential values. Provider
administrator credentials are forwarded only for the current preflight,
preparation, deployment, or destruction request.

Application authentication and cloud-resource access are separate concerns.
The deployed graph may still require AWS IAM Identity Center, Microsoft Entra,
GCP Identity-Aware Proxy, OIDC federation, or SAML-backed provider surfaces.
These mechanisms protect deployed cloud resources; they are not login methods
for Twin2MultiCloud itself.

\section{Management API and Flutter Client}
\label{sec:orchestrator}

\subsection{Management API Responsibilities}
\label{sec:management-api}

The Management API is the authoritative coordinator and persistence boundary.
Its central records are the owner profile, Digital Twin, CloudConnection,
CalculationRun, architecture selection, deployment manifest, deployment
operation, progress events, readiness evidence, access bundle, and cleanup
evidence. Provider-specific resource creation remains in the Deployer.

A draft Twin stores typed workload, device, event, simulator, state-machine,
and bounded function configuration. The primary extension slot is a validated
Python 3.11 telemetry processor. Its archive layout, entry point, dependencies,
size, declared capabilities, secret policy, and configuration are checked
against the versioned user-function contract. This preserves meaningful
university experimentation without accepting arbitrary executable project
trees.

Twin export and import use a versioned \texttt{.twin.zip} interchange format.
The archive contains allowlisted configuration and validated function sources;
it excludes CloudConnections, decrypted secrets, Terraform state, and
deployment outputs. Duplicate and import create a new draft with a new name.
They never mutate or implicitly destroy the source Twin.

\subsection{Twin State and Immutability}
\label{sec:twin-states}

The lifecycle uses the states \texttt{DRAFT}, \texttt{CONFIGURED},
\texttt{DEPLOYING}, \texttt{DEPLOYED}, \texttt{DESTROYING},
\texttt{DESTROYED}, \texttt{ERROR}, and \texttt{INACTIVE}. Configuration may be
edited in draft state. Once deployment starts, the calculation and deployment
inputs are immutable.

A changed deployed Twin is represented by a duplicated or imported draft and
receives a new calculation. The PoC does not perform in-place Terraform
updates, migrations, rollbacks, or continuous re-optimization. The same Twin
may be deployed again only after its previous infrastructure was successfully
destroyed.

\subsection{Durable Operations and SSE}
\label{sec:sse-streaming}

Deploy and Destroy can incur cost. Management therefore assigns an
idempotency key and correlation identifier to each operation, persists the
terminal state, and stores ordered progress events. Server-Sent Events (SSE)
provide live presentation, while \texttt{Last-Event-ID} supports reconnect and
bounded replay. SSE is only a transport optimization: the persisted operation
record remains authoritative, and a reconnect cannot issue another provider
mutation.

\subsection{Flutter Workflow}
\label{sec:flutter-ui}

The Flutter client presents the research workflow rather than a general
administration console. Its responsibilities are Twin configuration,
cost-result review, CloudConnection selection, readiness and repair,
deployment confirmation, operation progress, access handoff, telemetry
verification, and explicit Destroy. It does not contain provider SDKs, expose
manual provider overrides, or embed a Grafana administration surface.

\section{Cost Optimizer}
\label{sec:optimizer}

\subsection{API Boundary}
\label{sec:optimizer-api}

The Optimizer exposes one strict Six-layer calculation request. It accepts an
immutable workload, the exact profile reference, validated extension bindings,
and three pinned provider-pricing catalog references. Unknown or retired fields
are rejected. The response returns the monetary winner, exact monthly total,
cost ledger, RTA, RDS, pricing references, and bounded rejection diagnostics.

The normal API offers no objective selection, strategy registry, weighted
score, or manual component assignment. A repository-only evaluation utility
can reproduce an already enumerated candidate for the supervised evaluation;
that selector is rejected by the public request model and is never used by the
normal application workflow.

\subsection{Calculation Pipeline}
\label{sec:calc-pipeline}

Calculation proceeds in the following order:

\begin{enumerate}
    \item validate the immutable workload and exact profile/catalog digests;
    \item build provider options for all eight logical components;
    \item enumerate deterministic assignments and enforce architectural
    constraints, including co-location of hot storage and visualization;
    \item resolve every component and edge against the provider implementation
    profiles and component catalog;
    \item build an RDS and reject candidates that cannot be materialized;
    \item calculate component, route, transfer, bridge, and shared-pool costs
    exactly once; and
    \item rank admissible candidates by monetary monthly total with a stable
    assignment key as deterministic tie-breaker.
\end{enumerate}

This order prevents a cheap but incomplete provider bundle from winning. A
missing implementation mapping, formula, price, identity exchange, package,
or cleanup contract is a functional rejection rather than a zero-cost item.

\subsection{Cost Model}
\label{sec:cost-formulas}

Let $A(w)$ be the set of functionally admissible assignments for workload
$w$. For assignment $a$, the model sums component owners $I$ and cross-component
edge or route owners $E$ using one frozen pricing snapshot $p$:

\begin{equation}
    C(a,w,p) = \sum_{i \in I(a)} C_i(w,p)
             + \sum_{e \in E(a)} C_e(w,p).
\end{equation}

The selected architecture is therefore

\begin{equation}
    a^* = \operatorname*{arg\,min}_{a \in A(w)} C(a,w,p).
\end{equation}

Each ledger item records its formula reference, normalized quantities, unit,
price evidence, allocation owner, and subtotal. Cross-cloud data transfer and
bridge compute are attributed to explicit route owners; they are not hidden in
an adjacent layer. Only monetary cost participates in scoring. Latency,
sustainability, resilience, and weighted multi-objective optimization remain
conceptual future work.

\section{Cloud Deployer}
\label{sec:deployer}

\subsection{API and Package Boundary}
\label{sec:deployer-api}

The Deployer accepts a bounded operation package derived from one Twin and one
immutable Manifest v4. It no longer owns arbitrary named projects, project
history, active-project selection, or freely structured ZIP imports. Runtime
sources are attached to allowlisted extension slots and are checked again at
the deployment boundary.

\subsection{Readiness, Preparation, and Repair}
\label{sec:deploy-credential-validation}

Preflight is derived from the resolved graph rather than a fixed global list.
For each selected provider it checks credential identity and scope, required
services and APIs, region availability, permissions, and readable quota or
policy state where supported. Non-mutating readiness is distinct from
persistent provider preparation.

The preparation plan may register graph-required Azure resource providers or
enable graph-required GCP APIs after an explicit confirmation. It may also
identify an AWS outbound-identity prerequisite and present a supervised manual
step. Billing recovery, quota approval, organization-policy changes,
tenant-wide consent, legal approval, and provider-account creation remain
external responsibilities.

After a failed preparation, the user can retry an idempotent supported action,
follow typed manual instructions, or replace the selected CloudConnection.
Readiness is rerun after repair. Account-level capabilities are recorded
separately and are not removed by Twin Destroy.

\subsection{Deployment Pipeline}
\label{sec:deploy-pipeline}

The cost-incurring pipeline is deliberately linear:

\begin{enumerate}
    \item validate Manifest, RTA, RDS, workload, catalog, package, and
    CloudConnection digests;
    \item resolve graph-derived readiness and require any preparation evidence;
    \item require the explicit Apply confirmation and acquire the per-Twin
    operation lock;
    \item generate only allowlisted Terraform variables and bounded runtime
    packages from the RDS;
    \item execute Terraform initialization, validation, plan, and apply;
    \item run only the provider operations declared by the specification;
    \item collect typed access surfaces and perform the defined verification
    probes; and
    \item persist terminal operation and Terraform-output evidence.
\end{enumerate}

Destroy has its own confirmation and idempotency boundary. It invokes the
recorded Terraform state and declared cleanup contracts, then records provider
inventory and any residual resource. A partial cleanup is a failed operation,
not a successful Destroy with a warning.

\subsection{Terraform and Provider Adapters}
\label{sec:provider-abstraction}
\label{sec:iac-organization}
\label{sec:hybrid-model}

One versioned Terraform root module contains the reviewed AWS, Azure, and GCP
resources. RDS component selections activate the required resources and bind
their exact dimensions; provider assignment is not reconstructed from UI
fields. Runtime packages are content-addressed, and image references are
digest-pinned where a container is required.

Terraform owns declarative infrastructure. Narrow provider adapters are used
only when the selected service exposes a necessary operation outside the
Terraform resource model, such as uploading a provider-specific Twin model or
validating a runtime endpoint. Those operations are declared by the RDS and
have explicit verification and cleanup behavior; they are not an open-ended
post-deployment hook system.

\subsection{Resolved Variable Projection}
\label{sec:tfvars-pipeline}

The tfvars generator projects an allowlist from the RDS: selected providers,
regions, exact implementation IDs, normalized capacity dimensions, artifact
paths, and graph-derived routes. Credential values are injected through the
operation environment and are not serialized into portable Twin archives.
Terraform variable validation and native mock-provider tests verify that the
projection remains compatible with the root module without requiring cloud
credentials.

\section{Provider Implementation Mapping}
\label{sec:resource-mapping}

Table~\ref{tab:resource-mapping} summarizes the primary services in the
versioned component catalog. Each entry represents a reviewed service bundle,
not necessarily one managed product.

\begin{table}[htbp]
    \centering
    \caption{Primary provider services for the executable Six-layer profile.}
    \label{tab:resource-mapping}
    \small
    \begin{tabularx}{\textwidth}{@{} l X X X @{}}
        \toprule
        \textbf{Responsibility} & \textbf{AWS} & \textbf{Azure} & \textbf{GCP} \\
        \midrule
        Ingestion & IoT Core, IoT Commands, Lambda adapters & IoT Hub, Service Bus, Functions adapters & BifroMQ on GKE, load balancer, Pub/Sub adapter \\
        Eventing & Kinesis, SNS/SQS FIFO, Lambda & Event Hubs, Service Bus, Functions & Pub/Sub, Cloud Run services/worker pools \\
        Processing & Lambda, Step Functions & Functions, Logic Apps & Cloud Run, Workflows \\
        Hot storage & DynamoDB raw/rollup & Cosmos DB raw/rollup & Firestore raw/rollup \\
        Cool storage & S3 Standard-IA & Blob Cool & Cloud Storage Nearline \\
        Archive & S3 Glacier Deep Archive & Blob Archive & Cloud Storage Archive \\
        Twin state & IoT TwinMaker & Azure Digital Twins & Cloud Run Twin API and Firestore \\
        Visualization & Amazon Managed Grafana & Azure Managed Grafana & Grafana OSS on GKE \\
        \bottomrule
    \end{tabularx}
\end{table}

GCP is therefore a complete provider-hosted implementation in the executable
profile. Its Twin and visualization responsibilities are composed from hosted
software and managed infrastructure rather than excluded because no direct
managed equivalent exists. Version, capacity, identity, persistence,
networking, operations, cleanup, and cost ownership are explicit parts of the
bundle.

Twin2MultiCloud does not administer these visualization systems after
deployment. The access bundle provides the URL, provider, authentication kind,
assigned identity or bounded service-local credential, role, and readiness
state. The user opens the provider-hosted surface externally. The verification
contract proves that the bounded raw-history query is reachable; it does not
turn the application into a dashboard editor.

\section{End-to-End Workflow}
\label{sec:e2e-workflow}

The supported journey is:

\begin{enumerate}
    \item create, import, or duplicate a uniquely named draft Twin;
    \item configure the fixed workload, devices, events, simulator,
    state-machine behavior, and validated processor function;
    \item select user-owned CloudConnections for every provider required by
    the prospective calculation and deployment;
    \item calculate and review the cost winner, exclusions, assumptions, and
    exact evidence references;
    \item freeze the calculation and derive the deployment manifest;
    \item run graph-derived readiness, review any preparation plan, confirm it,
    and repair or replace a connection if necessary;
    \item review and explicitly confirm Apply;
    \item follow the durable operation through SSE or replay;
    \item open the returned access surface and execute one telemetry
    roundtrip; and
    \item explicitly confirm Destroy and reconcile residual resources.
\end{enumerate}

This workflow has no application login step and no manual layer override. The
six directed placements used by the thesis evaluation are materialized only by
the repository evaluation utility; the normal UI continues to deploy the
cost-selected candidate.

\section{Runtime Data Flow}
\label{sec:layer-data-flow}

\subsection{Layer Responsibilities}
\label{sec:telemetry-flow}

Ingestion accepts device telemetry and command outcomes. Eventing provides the
ordered, replayable domain-event boundary between ingestion, processing, hot
storage, and command delivery. Processing executes the validated user function
and emits canonical outcomes. Hot storage retains raw and bounded rollup data;
scheduled jobs move eligible records to cool and archive storage. Twin state
receives selected materializations and relationship changes. Visualization
queries bounded raw history from the co-located hot-storage provider.

The canonical event envelope is provider independent. Provider adapters map
native triggers and acknowledgements to this envelope, preserving the event
identifier, type, source, subject, time, schema version, correlation, and
payload. Delivery, retry, replay, terminal-failure, and observability bounds
are part of the versioned Eventing contract rather than undocumented runtime
behavior.

\subsection{Cross-Cloud Contracts}
\label{sec:cross-cloud-gluecode}

Only graph edges whose endpoints use different providers create cross-cloud
routes. Four edge contracts define the data semantics:

\begin{itemize}
    \item \texttt{canonical-domain-event.v1} for Eventing-mediated telemetry,
    processing outcomes, notifications, and device commands;
    \item \texttt{storage\_transition.v1} for hot-to-cool and cool-to-archive
    movement;
    \item \texttt{twin\_projection.v1} for selected state and relationship
    updates; and
    \item \texttt{raw\_history\_query.v1} for bounded visualization reads.
\end{itemize}

Raw-history queries and both storage transitions remain provider-local because
the executable PoC requires hot, cool, archive, and visualization to form one
co-located bundle. This constraint is enforced during candidate enumeration.
Canonical domain events and Twin projections have explicit directed
AWS/Azure/GCP route implementations with source-owned transfer, bridge,
identity, retry, and cost attribution. Cross-cloud storage migration remains
future work rather than an unverified deployment claim.

Cross-provider authentication uses workload identities and narrowly scoped
trust rather than embedding long-lived administrator credentials in runtime
functions. The exact mechanism depends on the directed pair and is declared by
the resolved graph. Administrator credentials are used to create these runtime
identities, not as application payloads.

\subsection{Failure and Recovery}
\label{sec:error-handling}

Failures are classified at their owning boundary: invalid intent, functional
incompatibility, pricing evidence failure, readiness blocker, preparation
failure, Terraform failure, runtime verification failure, or cleanup failure.
Public errors use stable codes and redacted diagnostics. Raw provider response
bodies and credentials are never returned to Flutter.

A failed Deploy does not authorize an automatic second Apply. The persisted
operation is inspected, and cleanup remains explicitly available. A failed
readiness or preparation step offers bounded repair before any Apply. A failed
Destroy retains residual evidence so the operator can complete cleanup without
the system claiming success.

\section{Evidence and State Management}
\label{sec:data-flow}

The Management database stores the state needed to reproduce and interpret an
experiment: typed Twin configuration, selected CloudConnection identifiers,
pricing references, workload and artifact digests, calculation result, RTA,
RDS, manifest, readiness and confirmation evidence, operation progress,
access handoff, verification result, Terraform outputs, Destroy result, and
residual classification.

Secrets and portable evidence have separate lifecycles. Encrypted provider
credentials remain in CloudConnections and are referenced by identifier.
Secret-free Twin archives and evaluation records can be shared. Deployment
administrator secrets, runtime-generated one-time credentials, and Terraform
state are excluded from exported research artifacts.

The offline gate validates schema synchronization, digest chains, formulas,
candidate admissibility, persistence, package projection, Terraform mock
plans, redaction, idempotency, replay, and cleanup contracts. These tests prove
deterministic implementation properties but are not presented as live-cloud
results. Live claims are reserved for the supervised, cost-capped evaluation
protocol described in Chapter~\ref{chap:evaluation}.
